\section{问题三：全向多源搜索模型的建立与求解}

\subsection{目标函数与状态表示}
设总移动距离、检测次数、频道切换次数、失败及成功清除次数分别为 $L,M,S,F,C$，则累计虚拟时间可核算为
\begin{equation}
  T=\frac{L}{5}+5M+S+3F+5C.
  \label{eq:virtual-time}
\end{equation}
优化目标是在清除全部存在源的约束下最小化 $T$。由于源数只知道在 $10$ 至 $16$ 之间，清除 $10$ 个源后不能直接退出。对每个频道维护
\begin{equation}
  \begin{gathered}
  \mathrm{UNKNOWN},\ \mathrm{FOUND},\ \mathrm{LOCALIZING},\\
  \mathrm{CLEARABLE},\ \mathrm{CLEARED},\ \mathrm{ABSENT\_CERTIFIED}.
  \end{gathered}
\end{equation}
六类状态，并记录方向观测、负观测、清除失败、可行域及覆盖完成度。

\subsection{无信号观测与全向覆盖}
对未清除的全向源，接收半径满足 $R_c\geq1000\,\mathrm{m}$。若在位置 $q$ 检测频道 $c$ 返回无信号，则有
\begin{equation}
  K_c\leftarrow K_c\setminus B(q,1000).
  \label{eq:omni-negative-update}
\end{equation}
只能排除半径 $1000\,\mathrm{m}$ 的圆盘，不能使用未知接收半径的上界 $1500\,\mathrm{m}$。若某频道在检测点集合 $P_c$ 上均返回无信号，且
\begin{equation}
  D_{1800}\subseteq\bigcup_{q\in P_c}B(q,1000),
  \label{eq:absence-certificate}
\end{equation}
则可确定该频道不存在未清除源，并将其标记为 $\mathrm{ABSENT\_CERTIFIED}$。

\subsection{中心加正六边形覆盖骨架}
取中心点 $p_{\mathrm{ctr}}=(0,0)$，并在半径 $a$ 的圆周上设置六个顶点
\begin{equation}
  p_k=a
  \begin{bmatrix}\cos(k\pi/3)\\ \sin(k\pi/3)\end{bmatrix},
  \qquad k=0,\ldots,5.
\end{equation}
中心检测覆盖 $r\leq1000\,\mathrm{m}$ 的区域。对 $1000\leq r\leq1800$，任一点与最近外圈顶点的方位差不超过 $30^\circ$，距离平方不超过
\begin{equation}
  f(r)=r^2+a^2-2ra\cos30^\circ.
\end{equation}
$f(r)$ 为凸函数，只需检查区间端点。外边界覆盖条件为
\begin{equation}
  1800^2+a^2-3600a\cos30^\circ\leq1000^2,
\end{equation}
其较小正根约为 $a_{\min}=1122.96\,\mathrm{m}$。取
\begin{equation}
  a=1135\,\mathrm{m}
\end{equation}
时，$r=1800\,\mathrm{m}$ 处的最坏距离约为 $994.81\,\mathrm{m}$，保留约 $5.19\,\mathrm{m}$ 的覆盖裕量。由中心依次访问六个外圈顶点且不返回中心，骨架路程为 $6a=6810\,\mathrm{m}$；相较 $a=1200\,\mathrm{m}$ 的方案减少 $390\,\mathrm{m}$，即节省 $78\,\mathrm{s}$ 移动时间。该布局以解析覆盖条件确定站点半径，使覆盖裕量与移动成本之间的取舍能够直接计算；其结论限于所构造的七站布局，不包含全局最优性。

\begin{figure}[H]
  \centering
  \includegraphics[width=0.67\textwidth]{F004-q3-coverage.pdf}
  \caption{全向源七站覆盖的几何证书；图中不表示实际运行轨迹}
  \label{fig:q3-coverage}
\end{figure}


\subsection{几何约束下的候选宏动作调度}
用公开请求与响应历史 $h_t$ 更新各频道的可行域、清除证书及覆盖进度，再生成候选集合 $\mathcal A(h_t)$。覆盖动作从最近的至多三个未访问站点中选择，串行检测未知频道及仍需补测的频道；已清除或已认证不存在的频道不再扫描。扫描优先使用当前频道，相邻站点交替频道顺序。

专门补测在可行域最小包围圆心的侧翼生成两个候选，方向垂直于最近示向度，偏移随包围圆半径限制在 $30$--$350\,\mathrm{m}$，每源至多专门补测三次。该实现使用问题二的交会几何思想，当前未调用其完整粗到细选点器。半径达到清除条件时前往包围圆心；返回近距离标志时优先原地清除。定位区域较小或补测次数达到阈值后，生成覆盖该区域的清除点序列，成功即停止。失败清除提供 $K_c\leftarrow K_c\setminus B(q,20)$ 的确定性信息。

候选包含覆盖、补测、保证清除、覆盖探测和退出五类宏动作。规则教师按估计成本排序，学习策略对同一集合评分，随后展开为串行请求。退出只在全部频道已清除或已认证不存在时开放；累计发现 $16$ 个不同频道后，可利用题面上界认证其余频道不存在。学习模块不能修改几何证书或放宽退出条件。

算法输入为公开历史 $h_t$ 与冻结参数，输出为一个合法宏动作。任务开始时，各频道均置为未知状态，机器人从题设初始位置生成首批覆盖候选；每执行一个宏动作便更新可行域和频道状态，再重新生成候选。若候选超出计算预算，则回退到规则教师；当全部频道取得清除或不存在证书时输出退出动作。这一接口把初始化、状态更新、回退与停止条件固定下来，参数化策略只改变候选的相对次序。

\begin{figure}[htbp]
\centering
\includegraphics[width=0.96\textwidth]{F001-controller-workflow.pdf}
\caption{在线控制器从公开历史更新证书、生成合法候选并执行后重新规划的闭环流程}
\end{figure}

\subsection{公开信息驱动的集合式候选评分}
覆盖结构决定哪些频道尚待搜索，几何证书决定哪些清除动作可靠，但两者并不唯一确定执行次序。为联合考虑信息获取和行动成本，将当前全局状态、各频道状态及候选动作分别编码。全局输入为 $10$ 维，频道输入为 $20\times9$ 维，每个候选为 $11$ 维，均由公开响应生成，真实源数与隐藏位置不输入策略。

频道编码经过两层四头注意力网络，隐藏维度为 $128$、前馈维度为 $256$，随后平均汇聚。候选编码与全局、频道汇聚编码拼接，输出动作分数 $\ell_\theta(h,a)$：
\begin{equation}
 \pi_\theta(a\mid h)=\frac{\exp\ell_\theta(h,a)}{\sum_{b\in\mathcal A(h)}\exp\ell_\theta(h,b)},\qquad a\in\mathcal A(h).
\end{equation}
变长候选以掩码处理；在线选择最高分宏动作，执行后重新生成候选。两问采用相同结构、分别冻结参数，每个网络含 $352386$ 个参数。频道聚合用于刻画待处理任务之间的联系，最终模型不包含候选与单频道之间的交叉注意力，也不包含粒子信念网络。

\subsection{以完整任务时间校准候选次序}
单独压缩移动路程可能推迟新源发现；只奖励成功清除又可能诱导过早停止。因此，将从任务开始到取得退出依据的全部虚拟时间纳入回报，采用
\begin{equation}
 r_t=-\frac{\Delta T_t}{100}-1000\mathbb I\{t\text{为失败终局}\},\qquad\gamma=1.
\end{equation}
其中失败终局包括未完成全部清除的回合。训练奖励不读取真实源数作归一化；逐局 $T/N$ 仅在评估与验证选模时计算。参数更新采用 PPO 剪切目标\cite{schulman2017}，记概率比为 $\rho_t=\pi_\theta(a_t\mid h_t)/\pi_{\rm old}(a_t\mid h_t)$，则策略项为
\begin{equation}
 \mathcal L_{\rm actor}=-\mathbb E\left[\min\left(\rho_t\widehat A_t,
 \operatorname{clip}(\rho_t,0.8,1.2)\widehat A_t\right)\right].
\end{equation}
优势估计采用 $\lambda=0.95$，并结合价值拟合与熵正则进行更新；每轮执行两个优化周期，近似 KL 超过 $0.03$ 时共同停止。训练时采样动作，部署时固定选择最高分动作。几何层始终独立维护合法候选及退出条件，策略无法通过提前退出换取较高回报。

方法的关键是把\textbf{发现覆盖、几何清除和不存在认证作为约束，把整局耗时作为调度目标}。PPO 是参数求解工具，其本身不作为本文原创算法；本文针对本题构造的覆盖证书、状态与动作接口共同限定了它的优化空间。前期相容场景搜索与数据聚合用于研究候选价值，相关实验保留在支撑材料，不将历史探索模块全部计入最终在线流程。

\subsection{完整性与独立实验结果}
在题设静态与有界误差条件下，七站覆盖保证每个存在的全向源至少可被一个站点检测；正观测持续更新保守可行域，必要时对该域构造有限清除覆盖。退出仅在 $20$ 个频道均已清除或已认证不存在时开放。因此，完整性来自覆盖与终止条件，不能用“连续若干次无信号”替代。

最终模型在 $3000$ 个独立研究场景中全部完成，共 $38961$ 个源全部发现并清除，\textbf{源级清除率与整局完成率均为 $100\%$}。含无源确认的平均总虚拟时间为 $3255.76\,\mathrm{s}$，逐局每源耗时的平均值为 $255.99\,\mathrm{s}/\text{源}$。相对同批父模型，平均总时间减少 $15.28\,\mathrm{s}$，每源耗时减少 $1.17\,\mathrm{s}$，配对区间见第九节。该结论是有限测试的实测表现，不把随机试验本身视为连续域完备性证明。

\subsection{正式测试记录的证据边界}
最终模型的上述结果来自自建研究环境。本次交付未提供问题三三次正式测试的完整成绩单，故不以研究结果填充正式成绩。提交时须由原始回放补齐案例编码、清除源数、平均定位清除时间及程序运行时间。
